\chapter{Reconstructing the paper's dataset and assumptions}\label{ch:data}
\paperloc{Section 2A, Tables 1--3, and the preprocessing in Section 2C.}

\section{The optical conditions that define this particular experiment}
An aberration-to-image map is conditional on many fixed inputs.
More explicitly, one should write
\[
\vct y=F(\vct c;\lambda,\NA,S,\mathcal M,\mathcal R,\mathcal Q),
\]
where $S$ is illumination, $\mathcal M$ is mask geometry and material
response, $\mathcal R$ is resist or contour extraction, and
$\mathcal Q$ is the set of sampled process conditions.
Changing any of these changes the map to be learned.

\begin{table}[htbp]\centering\small
\caption{The paper's Table 1, expressed with explicit meanings.
Values are reported from the publication; qualifications below are this
guide's interpretation.}
\begin{tabularx}{\textwidth}{@{}p{.31\textwidth}Y@{}}
\toprule Quantity & Reported value\\ \midrule
Vacuum wavelength & 13.5 nm\\
Projection NA & 0.33\\
Chief-ray angle at object & $6^\circ$\\
Illumination & Quasar: inner radius 0.6, outer radius 0.8,
angle $30^\circ$, rotation $45^\circ$\\
Absorbing topography & $\mathrm{Ta}_6\mathrm{N}_4$,
15 nm in wafer scale\\
Reflecting multilayer & Mo/Si, respectively 0.725 nm and 1 nm
in wafer scale\\
Focus settings & $-15,0,+15\nm$\\
Exposure-dose variation & $\pm5\%$\\
Zernike indices & $Z_2$ through $Z_{25}$\\
Coefficient sweep & $-20$ to $+20\,\mathrm{m}\lambda$,
step $1\,\mathrm{m}\lambda$\\
\bottomrule
\end{tabularx}
\end{table}

The optical constant database, complete multilayer termination,
number of periods, exact source discretization, polarization,
mask sidewalls, pupil apodization, and resist settings are not specified
well enough for a numerical reconstruction. A student should understand
their role rather than infer them from typical values.

\section{The intended pattern families in Table 2}
\begin{table}[htbp]\centering\small
\caption{Pattern specification reported in the paper's Table 2.
Each pitch list contains seven entries; rectangular pitch pairs are
read position by position.}
\begin{tabularx}{\textwidth}{@{}L{.22\textwidth}L{.14\textwidth}Y Y@{}}
\toprule Family & $(\mathrm{CD}_x,\mathrm{CD}_y)$ (nm)
& $p_x$ (nm) & $p_y$ (nm)\\ \midrule
Square vias & $(14,14)$ & 38, 40, 42, 46, 50, 80, 120
& 38, 40, 42, 46, 50, 80, 120\\
Rectangular vias & $(14,28)$ & 56, 60, 64, 70, 80, 100, 150
& 70, 74, 78, 84, 94, 114, 164\\
\bottomrule
\end{tabularx}
\end{table}
Through-pitch testing means holding a nominal feature type or size while
varying its neighbourhood spacing. This probes proximity effects and
different pupil sampling. It does not mean changing the wavelength or
continuously scanning every possible pitch.

The count suggested by two families and seven variants is 14 patterns.
Each contributes six scalar indicators:
\[
(\mathrm{CD}_x,\mathrm{CD}_y,\mathrm{PS}_x,\mathrm{PS}_y,
\mathrm{PVB}_x,\mathrm{PVB}_y).
\]
Thus $14\times6=84$ explains the reported network input width.
This is a consistent reconstruction of the feature count, not a
published list of the exact vector ordering.

\section{The conflicting plot-label table}
The paper's Table 3 assigns the following parameters to the bars labelled
REC and SQ. It is important to preserve this discrepancy:
\begin{table}[htbp]\centering\small
\caption{Paper Table 3 as printed. Every row uses CD $(14,28)$ nm for
REC and $(14,14)$ nm for SQ. These pitch assignments disagree with Table 2.}
\begin{tabular}{@{}crrrr@{}}
\toprule & \multicolumn{2}{c}{REC pitch (nm)}
& \multicolumn{2}{c}{SQ pitch (nm)}\\
Pattern & $p_x$ & $p_y$ & $p_x$ & $p_y$\\ \midrule
1&38&38&56&56\\
2&40&40&60&74\\
3&42&42&64&78\\
4&46&46&70&84\\
5&50&50&80&94\\
6&80&80&100&114\\
7&120&120&150&164\\
\bottomrule
\end{tabular}
\end{table}
The pitch families appear exchanged relative to the shape assignments
in Table 2. The first SQ pair $(56,56)$ also differs from the rectangular
pair $(56,70)$ in Table 2.
The text says REC means rectangle and SQ means square, but does not
resolve which table reflects the actual simulation layout.

When reading the graphs, retain the authors' labels and pattern numbers.
When performing this book's educational calculation, use explicitly
declared geometry. Do not claim that a chosen correction reproduces
the authors' dataset.

\section{Three groups of aberration inputs}
The stated simulation sets are:
\begin{enumerate}
\item \textbf{Individual-mode sweeps:} 41 coefficient values
$-20,-19,\ldots,+20\,\mathrm{m}\lambda$ for each of 24 indices,
with all others zero. This gives $41\times24=984$ rows.
\item \textbf{Definitive screening design:} 77 combined cases using
levels $-20,0,+20\,\mathrm{m}\lambda$.
\item \textbf{Random combinations:} 972 additional cases.
\end{enumerate}
Their nominal total is $984+77+972=2033$ before any filtering.
The paper does not give a complete split into training, validation,
and testing, nor the exact random distribution, seed, or retained row count.

The individual sweeps include the all-zero vector 24 times.
They therefore contain at most $984-23=961$ distinct vectors.
Duplicates between groups may further reduce the distinct count.
This matters in a replication: randomly splitting duplicate optical
cases between training and test sets would weaken the independence
of the test. The paper does not establish that such leakage occurred;
it simply leaves the handling undocumented.

\section{Why a screening design is used}
A local quadratic response in $K$ variables has
\[
1+K+K+\frac{K(K-1)}2
\]
coefficients: intercept, linear terms, squares, and pairwise interactions.
For $K=24$ this is $325$.
An efficient screening design seeks useful information about dominant
effects from far fewer than a full factorial set of runs, usually under
sparsity or low-order assumptions.

Three levels allow curvature to be distinguished from a simple straight
line. Carefully chosen combinations can separate some effects better
than arbitrary combinations. However, 77 runs cannot independently
estimate all 325 coefficients of an unrestricted quadratic model.
The exact design matrix and its aliasing properties are not supplied.
No specific orthogonality guarantee for these 77 rows can be reconstructed.

If every coefficient took all 41 grid values, a full 24-dimensional
grid would contain
\[
41^{24}=5.091110945\ldots\times10^{38}
\]
points. This count is a hypothetical grid comparison, not a claim that
the random group follows that grid. It illustrates why a few thousand
cases can only sample a minute portion of the possible coefficient space.

\section{What is one machine-learning sample?}
The natural forward row is
\[
\vct c^{(r)}
\stackrel{F}{\longrightarrow}
\vct y^{(r)}\in\R^{84},
\]
where the vector holds all patterns and metrics for aberration case $r$.
The inverse training row reverses the pair:
$\vct y^{(r)}$ is input, $\vct c^{(r)}$ is target.
The 14 patterns are therefore features within one optical case.
They are not automatically 14 independent realizations of the same input.

Keeping them separate is physically useful. For two patterns with
responses $\delta y_1=c$ and $\delta y_2=-c$, averaging produces zero
and loses all sensitivity to $c$. A feature vector $(c,-c)$ preserves it.
This is the paper's motivation for avoiding averages over pattern
performance during learning.

\section{The 24-versus-25 output issue}
The inclusive range $Z_2,\ldots,Z_{25}$ contains 24 coefficients.
The $984=41\times24$ count and plotted axes support 24 varied terms.
However, the prose repeatedly refers to the first 25 terms and the
reported network has 25 output neurons.

One possible explanation is inclusion of a fixed piston coordinate,
but the paper does not confirm that explanation. If piston is included
as an unconstrained output, intensity indicators cannot determine it.
A reconstruction must establish the exact output dictionary rather
than fabricate a twenty-fifth observable aberration.

\section{Filtering and scaling are part of the model}
Section 2C describes filtering poorly fitted or low-valued data and
mentions a 0.001 threshold, but its wording mixes input data and
output Zernike terms. It does not establish which rows or coordinates
were removed, the units before filtering, or whether the decision used
test information.

This ambiguity affects reproducibility and performance interpretation.
For example, removing all near-zero shift examples can improve a
percentage-error score while making predictions near the desired zero
shift less well supported. Filtering must be defined from the training
workflow and recorded; it is not just a harmless plotting choice.

\chapter{The inverse problem: existence, ambiguity, and feasible targets}\label{ch:inverse}
\section{Three separate questions}
Given desired metrics $\vct y_*$, ask:
\begin{enumerate}
\item \textbf{Existence:} is there any allowed wavefront with
$F(\vct c)=\vct y_*$, or sufficiently close?
\item \textbf{Uniqueness:} if a solution exists, is it the only one?
\item \textbf{Stability:} do small changes in desired metrics lead to
small changes in the required coefficients?
\end{enumerate}
A forward simulator may be perfectly deterministic even when the
inverse fails one or more of these tests.

\section{More indicators than coefficients does not settle uniqueness}
With about 24 varied coefficients and 84 metrics, there are more
reported numbers than unknowns. This does not imply either uniqueness
or nonuniqueness by itself. Metrics can be correlated, insensitive,
or redundant; nonlinear maps can also have multiple branches.

Locally, $J$ has 84 rows and 24 columns. If it has full column rank,
small coefficient changes are locally distinguishable by the measured
metrics. If its rank is lower, some first-order changes are invisible.
Even full local rank does not rule out two distant coefficient vectors
with the same metrics.

Piston is an exact unobservable direction if included. A simple
additional teaching example is $F(c)=c^2$: both $c=1$ and $c=-1$
give the same output. In symmetric optical settings, sign ambiguities
can arise from similar even dependence, but they should be established
for the actual source and mask rather than assumed from a mode name.

\section{The singular-value decomposition explains sensitivity}
Any real Jacobian can be written
\begin{equation}
J=U\Sigma V^T.
\label{eq:svd}
\end{equation}
The columns of $V$ are orthogonal coefficient-space directions.
Applying $J$ to direction $\vct v_j$ gives
$J\vct v_j=\sigma_j\vct u_j$.
The singular value $\sigma_j$ measures how much that coefficient
combination changes the measured outputs.

A small $\sigma_j$ means weak observability. Formally inverting this
component divides by $\sigma_j$, amplifying measurement and modelling
error. A zero value is a linear null direction.
This is a more precise concept than calling one single Zernike term
unimportant: a weak direction can be a correlated combination of many terms.

Before calculating an SVD, scale coordinates appropriately. Otherwise
replacing nm by milliwaves changes numerical singular values simply by
changing units. A useful scaled Jacobian uses coefficient allowance
scales in its columns and metric tolerance scales in its rows.

\section{Linear inverse and regularization}
For desired change $\delta\vct y$, a weighted least-squares inverse solves
\begin{equation}
\min_{\delta\vct c}
\frac12\norm{Q_y^{1/2}(J\delta\vct c-\delta\vct y)}^2
+\frac{\alpha}{2}\norm{L\delta\vct c}^2.
\label{eq:regularized}
\end{equation}
$Q_y$ weights metric errors; $L$ can penalize coefficient magnitude
or difficult-to-manufacture combinations; $\alpha\geq0$ balances the terms.
Differentiation gives
\begin{equation}
(J^TQ_yJ+\alpha L^TL)\delta\vct c=J^TQ_y\delta\vct y.
\end{equation}
With $Q_y=I$ and $L=I$, an SVD component is multiplied by
$\sigma_j/(\sigma_j^2+\alpha)$ rather than $1/\sigma_j$.
Small, unreliable directions are suppressed.
The penalty is a design preference, not additional optical data.

\section{Feasible outputs occupy a structured subset}
If $\vct c$ varies in a $K$-dimensional region, its smooth image in
84-dimensional metric space has at most $K$ local degrees of freedom.
Many metric combinations are therefore unattainable.

Being inside the minimum and maximum of each training feature is a
necessary range check, but not a sufficient physical-feasibility check.
For example, data satisfying $y_2=y_1^2$ with $-1\leq y_1\leq1$ have
coordinate ranges $y_1\in[-1,1]$ and $y_2\in[0,1]$.
The point $(0,1)$ lies inside both ranges but not on the curve.

The paper's 118 target vectors come from other simulations under the same
lithographic conditions and were not used in training or validation,
according to Section 3. That is stronger than choosing 84 coordinates
independently within a box. The paper nevertheless does not release the
vectors or their generation procedure, so their precise coverage is unknown.

\section{Why averaging inverse solutions can fail}
Suppose training pairs have identical input $y=1$ and two valid targets
$c=+1$ and $c=-1$ under $F(c)=c^2$.
A deterministic model trained by squared coefficient error minimizes
\[
L(\hat c)=\tfrac12[(\hat c-1)^2+(\hat c+1)^2]=\hat c^2+1.
\]
Its minimizer is $\hat c=0$, but $F(0)=0\ne1$.
The mean of valid inverse solutions is not necessarily valid.

More generally, squared-error regression favors the conditional mean
$\E[\vct c\mid\vct y]$, which can fall between distinct feasible branches.
A good test is therefore the forward consistency
\begin{equation}
\vct r_{\rm image}=F(G(\vct y))-\vct y,
\label{eq:forwardconsistency}
\end{equation}
not only the coefficient residual $G(\vct y)-\vct c_{\rm original}$.
The paper's forward re-simulation is an important part of its logic.
It checks whether a proposed alternative actually retains useful imaging.

\section{One output per input versus a distribution of solutions}
A deterministic feedforward network returns one coefficient vector
for one specified metric vector. Giving it 118 different target vectors
can produce 118 candidates. That does not mean the same trained network
samples all allowable wavefronts for a single fixed target.

To obtain diverse alternatives for the same target one could use multiple
optimization initializations, explicit latent inputs, or a conditional
generative model, followed by forward verification. These are possible
extensions, not methods reported in this paper.

\section{The feasible set is the real object of a budget}
Let permitted coefficient conditions be $\vct c\in\mathcal C$ and
imaging limits be $\vct y\in\mathcal Y_{\rm allowed}$.
The acceptable wavefront set is
\begin{equation}
\mathcal A=\{\vct c\in\mathcal C:F(\vct c)\in\mathcal Y_{\rm allowed}\}.
\label{eq:feasibleset}
\end{equation}
The paper identifies a finite collection of candidates in or near this set,
as judged by its simulator. The shape, volume, boundary, robustness, and
manufacturability of the entire set require further analysis.
Chapter \ref{ch:budget} develops the distinction between candidate
examples and a guaranteed tolerance region.
